% !TEX root = ../main.tex
% Modelo especificado en F00 §6. Formulación y comprobaciones desarrolladas aquí.
\chapter{Modelos Previos}
\label{ch:modelos-previos}

Antes de evolucionar un campo sobre Schwarzschild necesitamos comprobar
que el procedimiento numérico reproduce un problema cuya solución conocemos.
El caso elegido es Klein--Gordon en espacio plano con datos de un solo modo
espacial. Conserva la dispersión masiva que nos interesa y permite separar
el error del método de las dificultades introducidas por la geometría curva.
En este capítulo formulamos el problema y construimos su referencia exacta;
la discretización y los cálculos se presentan en el
capítulo~\ref{ch:simulacion}.

\section{Problema de referencia en espacio plano}

Consideramos un campo real $\Phi(x,t)$, con $c>0$ y
$\mu=m_\Phi c/\hbar>0$ constantes. Multiplicar
\eqref{eq:kg-plana-1d} por $c^2$ da
\begin{equation}
\Phi_{tt}-c^2\Phi_{xx}+c^2\mu^2\Phi=0.
\label{eq:kg-previo}
\end{equation}
La ecuación es de segundo orden temporal. Por eso no basta conocer el campo
en $t=0$: también se necesita su velocidad inicial. La variable espacial
requiere, a su vez, un dominio y condiciones en sus extremos. Son partes
distintas del problema: los datos iniciales seleccionan el estado del campo
y las condiciones de contorno especifican la geometría espacial de la prueba.

\section{Dominio y condiciones del problema}

Elegimos un intervalo $[0,L]$, con $L>0$, e identificamos sus extremos
mediante periodicidad. Para soluciones clásicas imponemos
\begin{equation}
\Phi(0,t)=\Phi(L,t),\qquad \Phi_x(0,t)=\Phi_x(L,t).
\label{eq:contorno-periodico}
\end{equation}
Así el campo y su primera derivada encajan al repetir el intervalo. Esta
elección admite perfiles armónicos exactos y elimina de esta prueba las
reflexiones de un borde artificial. El espacio modelado es periódico;
no equivale a imponer una condición de radiación en un espacio infinito.
Tampoco se interpreta como el contorno físico de un agujero negro.

Un perfil $e^{ikx}$ es periódico si $e^{ik(x+L)}=e^{ikx}$. Dividir por
$e^{ikx}\ne0$ exige $e^{ikL}=1$, es decir, $kL=2\pi j$ con $j\in\mathbb Z$.
Los números de onda permitidos son, por tanto,
\begin{equation}
k_j=\frac{2\pi j}{L}.
\label{eq:modos-periodicos}
\end{equation}
Su parte real, $\cos(k_jx)$, también es periódica. Para comprobar la
derivada basta observar que $\partial_x\cos(k_jx)=-k_j\sin(k_jx)$
y que $\sin(0)=\sin(2\pi j)=0$.

Seleccionamos uno de esos modos y prescribimos
\begin{equation}
\Phi(x,0)=A\cos(k_jx),\qquad \Phi_t(x,0)=0.
\label{eq:datos-previo}
\end{equation}
La velocidad inicial cero significa que el campo se libera desde ese perfil,
no que vaya a permanecer estático: la EDP determina su aceleración.
La amplitud $A$ tiene las unidades del campo. Para una prueba de frecuencia
se toma $A\ne0$; si $A=0$, la solución de referencia es nula. El modo $j=0$
es homogéneo en el espacio, mientras que para $j\ne0$ la longitud de onda
positiva es $L/|j|$. Los índices $j$ y $-j$ dan el mismo perfil de coseno.

\section{Solución analítica de referencia}

\subsection{Reducción del modo a un oscilador}
La función $\cos(k_jx)$ es autofunción de la segunda derivada espacial:
$\partial_x^2\cos(k_jx)=-k_j^2\cos(k_jx)$. Por ello buscamos una solución
que conserve ese perfil y cambie su amplitud con el tiempo,
\[
\Phi(x,t)=a_j(t)\cos(k_jx).
\]
Esta elección está motivada por el operador de coeficientes constantes y
el contorno periódico. Al sustituir en \eqref{eq:kg-previo},
\[
\bigl[a_j''(t)+c^2(k_j^2+\mu^2)a_j(t)\bigr]\cos(k_jx)=0.
\]
Como el coseno no es idénticamente cero y la igualdad vale para todo $x$,
el coeficiente temporal debe anularse:
\[
a_j''+\omega_{k_j}^2a_j=0,\qquad
\omega_{k_j}=c\sqrt{k_j^2+\mu^2}>0.
\]
La EDP se ha reducido a la misma estructura de oscilador que resolvimos
en Preliminares. Sus datos son $a_j(0)=A$ y $a_j'(0)=0$. Aplicando
\eqref{eq:oscilador-datos}, con $y_0$ reemplazado por $A$ y $v_0=0$,
obtenemos $a_j(t)=A\cos(\omega_{k_j}t)$. La solución candidata es
\begin{equation}
\Phi_{\mathrm{ex}}(x,t)=A\cos(k_jx)\cos(\omega_{k_j}t).
\label{eq:solucion-exacta-previo}
\end{equation}
El perfil es una onda estacionaria. La identidad del producto de cosenos
permite reconocerla como superposición de dos modos viajeros:
\[
\Phi_{\mathrm{ex}}(x,t)=\frac A2
\bigl[\cos(k_jx-\omega_{k_j}t)+\cos(k_jx+\omega_{k_j}t)\bigr].
\]
Ambos términos tienen la misma frecuencia y sentidos de fase opuestos;
su suma tiene velocidad inicial nula.

\subsection{Comprobación del problema completo}
Para comprobar la EDP calculamos
\[
\Phi_{\mathrm{ex},tt}=-\omega_{k_j}^2\Phi_{\mathrm{ex}},
\qquad \Phi_{\mathrm{ex},xx}=-k_j^2\Phi_{\mathrm{ex}}.
\]
Su residuo en \eqref{eq:kg-previo} es
\[
\bigl[-\omega_{k_j}^2+c^2k_j^2+c^2\mu^2\bigr]\Phi_{\mathrm{ex}}=0.
\]
En el instante inicial, $\cos(0)=1$ devuelve $A\cos(k_jx)$ y
\[
\Phi_{\mathrm{ex},t}=-A\omega_{k_j}\cos(k_jx)\sin(\omega_{k_j}t)
\]
se anula en cero. La periodicidad espacial se mantiene para todo $t$
porque el factor temporal es independiente de $x$. Se satisfacen, por
separado, la EDP, los dos datos iniciales y las dos condiciones de contorno.

\subsection{Unicidad y control de perturbaciones}
La verificación anterior establece la existencia de una solución para los
datos elegidos. Para justificar su uso como referencia debemos descartar
que otra solución clásica del mismo problema pueda dar una evolución
diferente. Si $w$ es la diferencia de dos soluciones reales $C^2$,
satisface la misma EDP y el mismo contorno periódico. Definimos
\[
\mathcal E_w(t)=\frac12\int_0^L
\left(c^{-2}w_t^2+w_x^2+\mu^2w^2\right)\,dx.
\]
Este funcional es no negativo. Derivarlo en tiempo da
\[
\mathcal E_w'(t)=\int_0^L
\left(c^{-2}w_tw_{tt}+w_xw_{xt}+\mu^2ww_t\right)\,dx.
\]
Integramos el segundo término por partes en $x$,
\[
\int_0^L w_xw_{xt}\,dx=[w_xw_t]_0^L-\int_0^L w_{xx}w_t\,dx,
\]
y agrupamos:
\[
\mathcal E_w'(t)=
\int_0^L w_t\left(c^{-2}w_{tt}-w_{xx}+\mu^2w\right)\,dx
+[w_xw_t]_0^L=0.
\]
La integral se anula por la EDP. El borde se cancela porque $w_x$ es
periódica y $w_t$ también lo es, al derivar en tiempo la igualdad
$w(0,t)=w(L,t)$. Así, $\mathcal E_w(t)=\mathcal E_w(0)$.
Si las dos soluciones tienen los mismos datos, $w(x,0)=w_t(x,0)=0$ y
$w_x(x,0)=0$; por tanto, $\mathcal E_w(0)=0$. Como $\mu>0$,
\[
\frac{\mu^2}{2}\int_0^L w(x,t)^2\,dx\le\mathcal E_w(t)=0.
\]
La continuidad implica $w=0$ y prueba la unicidad en esta clase.

El mismo cálculo controla perturbaciones de los datos. Definimos la norma
\[
\|f\|_{L^2(0,L)}=\left(\int_0^L|f(x)|^2\,dx\right)^{1/2}.
\]
La conservación anterior implica
\[
\|w(t)\|_{L^2}\le\frac{\sqrt{2\mathcal E_w(0)}}\mu,
\qquad
\|w_t(t)\|_{L^2}\le c\sqrt{2\mathcal E_w(0)}.
\]
Así, datos cercanos en la norma asociada a $\mathcal E_w$ producen campos
cercanos mientras existan las soluciones clásicas. Esta propiedad del
problema continuo será la referencia al analizar la estabilidad numérica;
no garantiza por sí sola que cualquier discretización sea estable.

\section{Objetivos de validación y parámetros}
\label{sec:parametros-previos}

La prueba numérica debe responder a dos cuestiones relacionadas pero distintas:
reproducir la solución \eqref{eq:solucion-exacta-previo} y recuperar su
frecuencia $\omega_{k_j}$. Comparar solo la forma de una gráfica no basta;
mediremos errores y su variación al refinar la malla. La elección de malla,
el integrador, las normas discretas y la extracción de frecuencia se
justifican en Simulación.

Las magnitudes físicas del modelo se recogen en la
Tabla~\ref{tab:parametros-previos}. Para ejecutar la validación se usan
las variables adimensionales de \eqref{eq:kg-adimensional}. En particular,
$\widehat L=L/L_0$, $\widehat A=A/\Phi_0$,
$\widehat k_j=k_jL_0$ y $\widehat\omega_j=\omega_{k_j}L_0/c$ satisfacen
\[
\widehat k_j=\frac{2\pi j}{\widehat L},\qquad
\widehat\omega_j=\sqrt{\widehat k_j^2+\widehat\mu^2}.
\]
Elegir $\widehat L=2\pi$, $\widehat\mu=0.7$ y $\widehat A=1$ proporciona
una familia de pruebas reproducibles, sin seleccionar una masa física ni
atribuir a estos valores significado observacional.

\begin{table}[htbp]
\centering\small
\begin{tabularx}{\textwidth}{@{}lXlX@{}}
\toprule
Símbolo & Magnitud & Unidad SI & Elección de validación \\
\midrule
$\Phi$, $A$ & Campo y amplitud & Como el campo & $\widehat A=1$ \\
$c$ & Velocidad de la luz & $\mathrm{m\,s^{-1}}$ & Absorbida en $\widehat t$ \\
$m_\Phi$ & Masa física del campo & kg & Sin masa SI seleccionada \\
$\mu$ & Parámetro $m_\Phi c/\hbar$ & $\mathrm{m^{-1}}$ & $\widehat\mu=0.7$ \\
$L$ & Longitud del dominio & m & $\widehat L=2\pi$ \\
$j$ & Índice del modo & $1$ & $1,2,3,4$ \\
$k_j$ & Número de onda permitido & $\mathrm{m^{-1}}$ & $\widehat k_j=j$ \\
$\omega_{k_j}$ & Frecuencia angular exacta & $\mathrm{s^{-1}}$ & $\widehat\omega_j=\sqrt{j^2+0.7^2}$ \\
\bottomrule
\end{tabularx}
\caption{Parámetros del modelo previo y valores adimensionales elegidos
para la primera validación. Los parámetros de discretización pertenecen
a Simulación.}
\label{tab:parametros-previos}
\end{table}
